% Marker macros
\providecommand{\BM}{\textbf{[B]}}
\providecommand{\KM}{\textbf{[K]}}
\providecommand{\SM}{\textbf{[S]}}
\providecommand{\SetzM}{\textbf{[ASSUMPTION]}}

\definecolor{ffgftblue}{RGB}{30,90,160}
\definecolor{proofgreen}{RGB}{0,110,50}
\newtcolorbox{kernbox}[1][]{breakable,enhanced,boxrule=0.6pt,arc=3pt,
  fonttitle=\bfseries\small,coltitle=white,top=4pt,bottom=4pt,left=6pt,right=6pt,
  colback=green!5,colframe=proofgreen,colbacktitle=proofgreen,#1}

\chapter{Mass-Bound Time Scales and Interferometry}
\label{chap:frequencyscales}

\begin{abstract}
From the FFGFT fundamental relation $\tilde{T}\cdot m=1$ follow
mass-bound Compton frequencies $f_m = mc^2/h$ as the only physically
grounded eigenfrequencies of a discrete substrate \KM. For all stable
particles these lie at least 13 orders of magnitude above the measurement
frequencies of current interferometers (MHz range). This yields a precise,
derivable statement about interferometry experiments for the detection of
space quantisation: not a design-induced blind spot but the mass-bound
nature of the time scale $\tilde{T}=1/m$ explains why Holometer-type
experiments at their current sensitivity cannot find a discreteness signal \KM.
A $K_\text{frak}$ correction to the Compton frequency of
$\Delta f_e/f_e = 100\xi_0 \approx 1.3\%$ is in principle spectroscopically
accessible, but not interferometrically \KM.
\end{abstract}

\section{Background: Two Approaches to Space Quantisation}

Experiments such as the Fermilab E-990 Holometer ($L=40\,\text{m}$,
$f_c = c/(2L) \approx 3.75\,\text{MHz}$) and the planned QUEST experiment
($L=1.84\,\text{m}$, $f_c \approx 81.5\,\text{MHz}$) search for a
quantisation signal of spacetime. The null results of these experiments
have been interpreted from various quarters as evidence for or against
specific discrete substrate models.

One concrete claim is that the interferometers are by design blind at
exactly the frequency at which a discrete D4 lattice with lattice constant
$a=\ell_P$ would resonate. This claim presupposes that the relevant
substrate frequency is $f \sim c/(2L)$.

The FFGFT framework provides a more precise answer to the question of
the relevant frequency scales --- from first principles, without parameter
fitting.

\section{FFGFT Eigenfrequencies from $\tilde{T}\cdot m=1$}

\subsection{The Mass-Bound Time Scale}

\SetzM\ In FFGFT, time is not an independent external quantity but is bound to
mass by the fundamental relation $\tilde{T}\cdot m=1$ (in natural units;
in SI: $\tilde{T}\cdot m = \hbar/c^2$).

\KM\ Every massive particle therefore carries an intrinsic time scale:
\begin{equation}
\tilde{T}_m = \frac{\hbar}{mc^2}, \qquad f_m = \frac{mc^2}{h}.
\label{eq:comptonfrequency}
\end{equation}
This is the Compton frequency --- not an additional postulate but a geometric
consequence of the $T^4$ structure and the duality relation.

\subsection{External Inputs}

\SetzM\ This document uses the following standard values as external inputs;
they are declared here and not derived:

\begin{center}
\begin{tabular}{lll}
\toprule
Quantity & Value & Source \\
\midrule
$m_e$   & $9.1094\times10^{-31}\,\text{kg}$ & PDG \\
$m_\mu$ & $1.8836\times10^{-28}\,\text{kg}$ & PDG \\
$m_p$   & $1.6726\times10^{-27}\,\text{kg}$ & PDG \\
$m_n$   & $1.6749\times10^{-27}\,\text{kg}$ & PDG \\
$h$     & $6.6261\times10^{-34}\,\text{J\,s}$ & SI (derived from $\xi$) \\
$c$     & $2.9979\times10^{8}\,\text{m/s}$  & SI (derived from $\xi$) \\
\bottomrule
\end{tabular}
\end{center}

\subsection{Numerical Values}

\KM

\begin{kernbox}[title={Compton frequencies of stable particles (SI)}]
\begin{center}
\begin{tabular}{lrr}
\toprule
Particle & $f_m = mc^2/h$ & $R_m = \hbar/(mc^2)$ \\
\midrule
Electron  & $1.2356\times10^{20}\,\text{Hz}$ & $1.288\times10^{-21}\,\text{m}$ \\
Muon      & $2.5548\times10^{22}\,\text{Hz}$ & $6.230\times10^{-24}\,\text{m}$ \\
Proton    & $2.2687\times10^{23}\,\text{Hz}$ & $7.015\times10^{-25}\,\text{m}$ \\
Neutron   & $2.2719\times10^{23}\,\text{Hz}$ & $7.006\times10^{-25}\,\text{m}$ \\
\bottomrule
\end{tabular}
\end{center}
\end{kernbox}

\subsection{Comparison with Interferometer Frequencies}

\KM
\begin{align}
\frac{f_e}{f_c^\text{Holometer}} &= \frac{1.2356\times10^{20}\,\text{Hz}}{3.747\times10^{6}\,\text{Hz}}
\approx 3.3\times10^{13}, \\
\frac{f_e}{f_c^\text{QUEST}} &\approx 1.5\times10^{12}.
\end{align}

The smallest mass-bound frequency (electron) thus lies 13 orders of
magnitude above the measurement frequencies of current interferometers.

\section{Lattice Dispersion Correction at the FFGFT Scale}

\KM\ For a discrete substrate with lattice constant $a$, the lattice dispersion
correction to linear dispersion is of order $(k\cdot a)^2$. At
$k_1 = \pi/L$ this gives:
\begin{equation}
\left.\frac{\delta\omega}{\omega}\right|_\text{lattice} \sim \left(\frac{\pi a}{L}\right)^2.
\end{equation}

\KM

\begin{center}
\begin{tabular}{lll}
\toprule
Lattice scale $a$ & Holometer ($L=40\,\text{m}$) & QUEST ($L=1.84\,\text{m}$) \\
\midrule
$\ell_P = 1.62\times10^{-35}\,\text{m}$ & $\sim 10^{-72}$ & $\sim 10^{-70}$ \\
$R_{m_e} = 1.29\times10^{-21}\,\text{m}$ & $\sim 10^{-44}$ & $\sim 10^{-42}$ \\
$R_{m_p} = 7.0\times10^{-25}\,\text{m}$ & $\sim 10^{-51}$ & $\sim 10^{-49}$ \\
\bottomrule
\end{tabular}
\end{center}

Even at the Compton scale of the electron --- which in the FFGFT framework
is the \emph{physical} fundamental scale, not $\ell_P$ --- the lattice
dispersion correction is $\sim10^{-44}$, far beyond any measurable
sensitivity.

\section{The $K_\text{frak}$ Correction and Its Measurability}

\KM\ The fractal correction $K_\text{frak} = 1 - 100\xi_0$ with $\xi_0 = 4/30000$
acts exclusively on absolute SI values, not on the pullback isometry itself
(Doc.~190, R96, R98). Applied to the Compton frequency of the electron:
\begin{equation}
\Delta f_e = f_e \cdot 100\xi_0 = f_e \cdot \frac{400}{30000}
\approx 1.65\times10^{18}\,\text{Hz}.
\end{equation}
This corresponds to a relative deviation of $100\xi_0 \approx 1.33\%$
from the Compton frequency. This correction is in principle spectroscopically
accessible (precision spectroscopy at $\sim10^{20}\,\text{Hz}$, i.e.\ hard
X-ray to gamma range), but not with a Michelson interferometer in the MHz
range.

\section{The Central Statement on Interferometry Experiments}

\begin{kernbox}[title={Conclusion: Why Holometer-type experiments find no signal}]
\KM\ In FFGFT the relevant frequency scale for space quantisation is not
$f \sim c/(2\ell_P)$ (universal Planck frequency) but $f_m = mc^2/h$
(mass-bound Compton frequency). This scale differs for every particle and
lies for all stable particles at least 13 orders of magnitude above the
measurement frequencies of current interferometers.

It follows directly: Holometer-type experiments find nothing because the
time scale of discreteness is not universal but mass-bound --- not because
a design-induced blind spot of the instrument lies exactly at the substrate
resonance frequency.

The statement is falsifiable: an interferometer signal at frequency $f$
would require a massless or ultralight field with $m = hf/c^2$. For
$f = 3.75\,\text{MHz}$ this gives $m \approx 2.8\times10^{-44}\,\text{kg}
\approx 3\times10^{-14}\,m_e$ --- no known particle.
\end{kernbox}

\section{Relation to the D4 Lattice Argument}

\KM\ The D4 lattice with lattice constant $a=\ell_P$ has 24 nearest neighbours
(8 of type $(\pm1,0,0,0)$ and 16 of type
$(\pm\tfrac{1}{2},\pm\tfrac{1}{2},\pm\tfrac{1}{2},\pm\tfrac{1}{2})$).
The dispersion relation of a scalar field on this lattice is:
\begin{equation}
\omega^2(\mathbf{k}) = \frac{c^2}{\ell_P^2}\sum_{\boldsymbol{\delta}}
\bigl(2-2\cos(\mathbf{k}\cdot\boldsymbol{\delta}\,\ell_P)\bigr).
\end{equation}
For $k_1 = \pi/L$ with $L=40\,\text{m}$ one has
$k_1\ell_P \approx 1.3\times10^{-36}$, hence $k_1\ell_P \ll 1$: the
dispersion relation is exactly linear in this regime, identical to the
continuum up to order $(k_1\ell_P)^2\sim10^{-72}$. No special lattice
resonance exists at $f_c = c/(2L)$ --- the free spectral range $c/(2L)$
is an optical property of the interferometer, not a lattice property.

FFGFT replaces the D4 lattice argument with a stronger statement: not the
lattice dispersion but the mass-bound time structure $\tilde{T}\cdot m=1$
determines at which frequency a signal is to be expected.

\section*{Notation}

\begin{tabular}{ll}
\textbf{[ASSUMPTION]} & Axiom or declared external input -- not derived \\
\textbf{[K]}          & Core -- derived from $\xi$ and the assumptions, numerically verified \\
\textbf{[B]}          & Bridge -- algebraically proved \\
\textbf{[S]}          & Sketch -- plausible, not yet fully worked out \\
\end{tabular}

Verification script: \texttt{pruef\_335\_paul\_frequenzen.py} (5/5 assertions).

\emph{Declared 27~August 2026}
